%% This is file `example.tex',
%% generated with the docstrip utility.
%% The original source files were:
%% coppe.dtx  (with options: `example')
%% This is a sample monograph which illustrates the use of `coppe' document
%% class and `coppe-unsrt' BibTeX style.
%% \CheckSum{1311}
%% \CharacterTable
%%  {Upper-case    \A\B\C\D\E\F\G\H\I\J\K\L\M\N\O\P\Q\R\S\T\U\V\W\X\Y\Z
%%   Lower-case    \a\b\c\d\e\f\g\h\i\j\k\l\m\n\o\p\q\r\s\t\u\v\w\x\y\z
%%   Digits        \0\1\2\3\4\5\6\7\8\9
%%   Exclamation   \!     Double quote  \"     Hash (number) \#
%%   Dollar        \$     Percent       \%     Ampersand     \&
%%   Acute accent  \'     Left paren    \(     Right paren   \)
%%   Asterisk      \*     Plus          \+     Comma         \,
%%   Minus         \-     Point         \.     Solidus       \/
%%   Colon         \:     Semicolon     \;     Less than     \<
%%   Equals        \=     Greater than  \>     Question mark \?
%%   Commercial at \@     Left bracket  \[     Backslash     \\
%%   Right bracket \]     Circumflex    \^     Underscore    \_
%%   Grave accent  \`     Left brace    \{     Vertical bar  \|
%%   Right brace   \}     Tilde         \~}
%

%Escolher uma das 3 possibilidades no primeiro argumento
% msc - dissertacao de mestrado
% dscexam - qualificacao
% dsc - tese de doutorado
% tcc - trabalho de conclusão de curso
%
\documentclass[tcc,divps,pdftex]{tcc_modelo}
\usepackage[brazil]{babel}
\usepackage[utf8]{inputenc}
\usepackage{amsmath,amsthm,amsfonts,amssymb,enumerate}
\usepackage{setspace}
\usepackage[lined,boxed,portugues]{}
\usepackage{booktabs}
\usepackage[portuguese,algoruled,longend,linesnumbered]{algorithm2e}
\usepackage{natbib}
\usepackage{subfigure}
\usepackage{subeqnarray}
\usepackage{enumerate}
\usepackage{multicol}
%\usepackage{slashbox}

\usepackage{cases}
\usepackage{float}
\usepackage{soul}

\usepackage{pdfpages}

\usepackage{float}

\usepackage{listings}

\usepackage{xcolor,colortbl}         % Pacote de cores
\usepackage{bm}


\renewcommand{\lstlistingname}{Algoritmo}
\renewcommand{\algorithmcfname}{Pseudocódigo}


\newcommand{\R}{\mathbb{R}}
\newcommand{\N}{\mathbb{N}}
\newcommand{\Z}{\mathbb{Z}}
\newcommand{\Q}{\mathbb{Q}}
\newcommand{\K}{\mathbb{K}}
\newcommand{\I}{\mathbb{I}}
\newcommand{\id}{\mathbf{1}}
\newcommand{\U}{\mathcal{U}}
\newcommand{\V}{{\cal V}}

\newcommand{\mat}[1]{\mbox{\boldmath{$#1$}}}


\newtheorem{teo}{Teorema}[chapter]

%\input{hifenizacao}



\sloppy
\hyphenpenalty=10000

%\makelosymbols
%\makeloabbreviations
%\makenomenclature
% \makeglossary
%\usepackage{titlesec}
%\titleformat{\section}{\bfseries\normalsize}{\uppercase{\sectiontitle}\thesection.\quad}{1em}{}

\begin{document}

\course{ENGCOMP} %engenharia Computacional
\institute{FACENG_ICE} %Faculdade de Engenharia e Instituto de Ciências Exatas

\title{Quantificação de Estratégias Discricionárias de Trading: Uma Abordagem Neuro-Fuzzy Aplicada aos Conceitos de Smart Money.}%Possível título
%\foreigntitle{To be defined...}

\author{}{Felipe da Costa Pinto Vilela}

\advisor{Prof.}{Luciana Conceição Dias}{Campos}{Dr.}{University}
\coadvisor{}{}{}{}{}
\coadvisorII{}{}{}{}{}

\examinerI{Prof.}{Membro da Banca 1}{D.Sc.}{Universidade }
\examinerII{Prof.}{Membro da Banca 2}{D.Sc.}{Universidade }
\examinerIII{}{}{}{}
\examinerIV{}{}{}{}



%definir data da defesa
\date{20}{07}{2026}

%Palavras chaves do resumo em portugues
\keyword{Palavra chave 1}
\keyword{Palavra chave 2}
\keyword{Palavra chave 3}
\keyword{Palavra chave 4}
\keyword{Palavra chave 5}

%Palavras chaves do abstract em ingles
\keywordEN{Keyword 1}
\keywordEN{Keyword 2}
\keywordEN{Keyword 3}
\keywordEN{Keyword 4}
\keywordEN{Keyword 5}


\inserecapa

\maketitle

\frontmatter

\dedication{A quem você dedica o trabalho...}

 \begin{agradecimentos}

Agradeço ...
 \end{agradecimentos}
 
 \epigrafe{``You only fail when you stop trying''  \flushright{(Thomas Edison)}}

\begin{abstract}
Resumo ...  
\end{abstract}

\begin{foreignabstract}
Abstract ...
\end{foreignabstract}



% \clearpage


\listoffigures  %Com titulo : lista de ilustracoes
\listoftables

%----Glossário ------------
\clearpage
\clearpage
\begin{center}
\textbf{LISTA DE ABREVIAÇÕES E SIGLAS}
\end{center}

\vspace*{1.cm}

\doublespacing  \begin{tabular}{l l}
CTH& 	 Células-Tronco Hematopoiéticas\\
FOU& 	 \textit{First Order Upwind}\\
LPS& 	 Lipopolissacárido\\
MAC& 	 Moléculas de Adesão Celular\\
MDF&	 Método dos Diferenças Finitas\\
MEC&     Matriz Extracelular\\
MEF&	 Método dos Elementos Finitos\\
MVF&	 Método dos Volumes Finitos\\
PFI&	 Pressão do Fluido Intersticial\\
RMC& 	 Ressonância Magnética Cardiovascular\\
SII&	 Sistema Imune Inato\\
SUPG&	 \textit{Streamline upwind} Petrov–Galerkin\\
\end{tabular}  \thispagestyle{empty}
 %arquivo abreviacoes.tex não é automatizado :/

\tableofcontents


\mainmatter

%------------------------------------------------------------------------------

%deixei uma introdução para tomar como exemplo, algumas formataçoes não foram automaticas e são feitas manualmente no .tex
 \chapter{INTRODUÇÃO}


\section{\uppercase{Motivação e Contextualização}}

O curso de Engenharia Computacional tem como principal característica a interdisciplinaridade e é uma proposta conjunta da Faculdade de Engenharia e do Instituto de Ciências Exatas nos mesmos moldes do Programa de Pós-Graduação em Modelagem Computacional.

A proximidade a este curso de pós-graduação garante ao aluno uma formação em um ambiente de pesquisa, de modo que este se habitue a questionar, buscar novas soluções, verificar suas ideias e compará-las com as de outros, o que certamente se constituirá em uma vantagem comparativa no mercado de trabalho.

Já no curso de Engenharia Computacional, o computador passa de objetivo central a meio. O objetivo do curso é o estudo de diferentes áreas da Ciência e das Engenharias através de ferramentas computacionais. O curso faz uma integração horizontal de diversas áreas do conhecimento: software e sistemas computacionais, matemática computacional, modelagem computacional, aplicadas à resolução de problemas diversos da Ciência e das Engenharias, como os da Engenharia Civil, Elétrica, Ambiental, de Produção, Mecânica, de Petróleo, Biologia, Física, Química, etc.

\section{\uppercase{Objetivos}}

O curso de Engenharia Computacional proporciona uma formação essencialmente multidisciplinar, visando a formulação, análise, implementação e aplicação de modelos matemáticos, métodos numéricos e sistemas computacionais para a solução de problemas de engenharia. Tal formação permite ao egresso do curso a inserção em um segmento do mercado de trabalho que está voltado aos recentes desenvolvimentos tecnológicos, além de também poder se constituir em um primeiro passo para uma carreira acadêmica multidisciplinar de origem. De uma forma geral, os egressos estarão capacitados a resolver problemas complexos que permeiam as áreas de computação e engenharia~\citep{lobosco2018design}.

Esta possibilidade de formação se encaixa em uma nova área de conhecimento que vem se consolidando nas últimas décadas, reunindo especializações nascidas em diferentes ramos da engenharia e da ciência. Internacionalmente, esta área de formação é denominada Computational Science and Engineering (Ciência e Engenharia Computacional).

Mais especificamente, o Engenheiro Computacional é um profissional com formação fortemente interdisciplinar, capaz de atuar na análise, modelagem e simulação de fenômenos físicos. O futuro profissional formado em Engenharia Computacional poderá, dependendo das escolhas a seu critério, ao longo de sua trajetória acadêmica, identificar, formular, modelar e desenvolver modelos matemáticos para a representação de problemas complexos nas mais distintas áreas da engenharia, dentre as quais podemos citar: elétrica, eletrônica, telecomunicações, agronomia, petróleo, mecânica, naval, civil e estrutural, aeronáutica, ambiental, materiais, e biomecânica, dentre outras. É também capaz de especificar e implementar a representação computacional dos modelos propostos, bem como supervisionar a manutenção e operação dos sistemas computacionais desenvolvidos. Coordena e supervisiona equipes de trabalho, realiza estudos de viabilidade técnico-econômica, executa e fiscaliza serviços técnicos em computação; e efetua perícias e avaliações, emitindo laudos e pareceres. Em suas atividades, considera a ética, a segurança, a legislação e os impactos socioambientais dos modelos e sistemas desenvolvidos.

Atendidos os conteúdos do núcleo básico da Engenharia, os conteúdos profissionalizantes do curso de Engenharia Computacional são abertos, dependendo das escolhas do aluno e das características do corpo docente. Entretanto aconselha-se um conteúdo profissionalizante fortemente interdisciplinar, que dê ao aluno uma complementação de sua formação matemática e de computação, englobando temas como modelagem computacional, métodos numéricos, métodos dos elementos finitos, método dos elementos de contorno, métodos estatísticos, linguagens de programação, engenharia de software, banco de dados, computação paralela e distribuída, análise e projeto de algoritmos. Também sugere-se que aja uma interação com outros cursos de Engenharia, permitindo ao aluno um conhecimento em áreas como elétrica, eletrônica, telecomunicações, agronomia, petróleo, mecânica, naval, estrutural, aeronáutica, ambiental, materiais, e biomecânica, dentre outras. O objetivo é permitir ao aluno se aprofundar na modelagem de problemas ligados as áreas das Engenharias que mais lhe atraiam, naturalmente observadas as restrições de disponibilidade de disciplinas. No caso específico de nosso curso, pela fato do Departamento de Mecânica Aplicada e Computacional ser um dos idealizadores do curso, se oferecem aos alunos uma maior especialização na área de modelagem de estruturas e de mecânica computacional. Contudo, deve-se destacar que a existência de outros cursos de graduação em Engenharia na UFJF abre a possibilidade para a criação de ênfases na modelagem de problemas em outras áreas da Engenharia, como Sanitária e Ambiental, Mecânica, Produção e Elétrica.

Como todos os cursos que tem ingresso através do curso de Ciências Exatas da UFJF, o aluno fará jus ao diploma referente ao Bacharelado Interdisciplinar em Ciências Exatas e, após a conclusão dos requisitos específicos, do diploma em Engenharia Computacional. O curso possibilita ainda, ao aluno que assim o desejar,  a obtenção de diplomas adicionais em outros cursos de segundo ciclo do  Bacharelado Interdisciplinar em Ciências Exatas, como Ciência da Computação, Matemática, Física, Química, Estatística, dentre outros.




% \input{./capitulos/cap2.tex}

% \input{./capitulos/cap3.tex}

% ...

% \input{./capitulos/cap_n.tex}


%------------------------------------------------------------------------------

%\renewcommand{\bibname}{REFERÊNCIAS}
\label{bib:begin}
\renewcommand\bibname{\centerline{REFERÊNCIAS}\global\def\bibname{REFERÊNCIAS}}

\bibliographystyle{bibformat}
\bibliography{referencias}
\label{bib:end}

%\appendix
%\input{./capitulos/appendix.tex}


\end{document}
